% Ch. 24 : promotion d'un artefact unique
\documentclass[border=6pt]{standalone}
\usepackage{coursfig}
\begin{document}
\begin{tikzpicture}[x=1mm,y=1mm]
  \node[box=Ocre, text width=62mm, minimum height=18mm] (img) at (0,30)
       {\ico[Ocre]{box}\enspace\textbf{Une seule image, construite une fois}\\[1mm]{\small\ttfamily listify-backend@sha256:4f2c\ldots}};
  \foreach \x/\e/\c/\cfg [count=\k] in {-74/dev/Enc/{1 réplique\\base jetable}, 0/préproduction/Sea/{2 répliques\\copie anonymisée}, 74/production/Plum/{3 répliques\\base réelle}}{
    \node[box=\c, text width=42mm, minimum height=24mm] (e\k) at (\x,-10) {\textbf{\e}\\[1mm]{\normalsize\color{Muted}\cfg}};
    \draw[flow=Ocre] (img.south) -- ++(0,-6) -| (e\k.north);
  }
  \draw[flow, line width=1pt] (e1) -- node[elabel, above=1mm, fill=none]{promotion} (e2);
  \draw[flow, line width=1pt] (e2) -- node[elabel, above=1mm, fill=none]{promotion} (e3);
  \node[note, anchor=north, text width=170mm, align=flush center] at (0,-26) {ce qui change d'un environnement à l'autre : la configuration (valeurs, secrets) ; ce qui ne change jamais : l'image, identifiée par son digest};
\end{tikzpicture}
\end{document}
